\documentclass[border=10pt]{standalone}
\usepackage{ctex}
\usepackage{amsmath,amssymb}
\usepackage{tikz}
\usepackage{tikz-feynman}
\tikzfeynmanset{compat=1.1.0}

\begin{document}
\begin{tikzpicture}

% ===== Title =====
\node[font=\Large\bfseries] at (8, 0) {Dyson 级数 $\longrightarrow$ Wick 收缩 $\longrightarrow$ Feynman 图};

% ===== n=0 =====
\node[font=\bfseries, blue] at (0, -2) {$n=0$};
\node at (2.5, -2) {$U^{(0)} = 1$};
\node[red] at (5, -2) {真空（无相互作用）};

% ===== n=1 =====
\node[font=\bfseries, blue] at (0, -4) {$n=1$};
\node at (3, -4) {$\displaystyle\frac{-i}{\hbar}\int dt_1\, T[H_I(t_1)]$};

% Feynman diagram: single vertex with 4 legs (phi^4)
\begin{feynman}
  \vertex (v1) at (8, -4);
  \vertex (a1) at (7, -3.2);
  \vertex (a2) at (7, -4.8);
  \vertex (a3) at (9, -3.2);
  \vertex (a4) at (9, -4.8);
  \diagram* {
    (a1) -- (v1),
    (a2) -- (v1),
    (a3) -- (v1),
    (a4) -- (v1),
  };
  \filldraw[black] (v1) circle (2pt);
\end{feynman}
\node[font=\small] at (8, -5.3) {单顶点（tadpole）};

% ===== n=2 =====
\node[font=\bfseries, blue] at (0, -7) {$n=2$};
\node at (3.2, -7) {$\displaystyle\frac{(-i)^2}{2!\,\hbar^2}\int dt_1 dt_2\, T[H_I(t_1)H_I(t_2)]$};

% Feynman diagram: two vertices connected by propagators
\begin{feynman}
  \vertex (v1) at (7, -7);
  \vertex (v2) at (9.5, -7);
  \vertex (a1) at (6, -6.2);
  \vertex (a2) at (6, -7.8);
  \vertex (a3) at (10.5, -6.2);
  \vertex (a4) at (10.5, -7.8);
  \diagram* {
    (a1) -- (v1),
    (a2) -- (v1),
    (v1) -- [edge label'=$G_F$] (v2),
    (v2) -- (a3),
    (v2) -- (a4),
  };
  \filldraw[black] (v1) circle (2pt);
  \filldraw[black] (v2) circle (2pt);
\end{feynman}
\node[font=\small] at (8.25, -8) {树图（2 顶点 + 1 内线）};

% Second topology for n=2: vacuum bubble
\begin{feynman}
  \vertex (v1b) at (12, -7);
  \vertex (v2b) at (14.5, -7);
  \diagram* {
    (v1b) -- [bend left=60, edge label=$G_F$] (v2b),
    (v1b) -- [bend right=60, edge label'=$G_F$] (v2b),
  };
  \filldraw[black] (v1b) circle (2pt);
  \filldraw[black] (v2b) circle (2pt);
\end{feynman}
\node[font=\small] at (13.25, -8) {真空泡（2 内线）};

% ===== n=3 =====
\node[font=\bfseries, blue] at (0, -10.5) {$n=3$};
\node at (3.5, -10.5) {$\displaystyle\frac{(-i)^3}{3!\,\hbar^3}\int dt_1 dt_2 dt_3\, T[H_I H_I H_I]$};

% Triangle vacuum diagram
\begin{feynman}
  \vertex (v1c) at (8, -9.8);
  \vertex (v2c) at (7, -11.2);
  \vertex (v3c) at (9, -11.2);
  \diagram* {
    (v1c) -- [edge label=$G_F$] (v2c),
    (v2c) -- [edge label'=$G_F$] (v3c),
    (v3c) -- [edge label=$G_F$] (v1c),
  };
  \filldraw[black] (v1c) circle (2pt);
  \filldraw[black] (v2c) circle (2pt);
  \filldraw[black] (v3c) circle (2pt);
\end{feynman}
\node[font=\small] at (8, -12) {三角形真空图};

% Sunset diagram
\begin{feynman}
  \vertex (v1d) at (11, -10.5);
  \vertex (v2d) at (13.5, -10.5);
  \diagram* {
    (v1d) -- [bend left=45, edge label=$G_F$] (v2d),
    (v1d) -- [edge label=$G_F$] (v2d),
    (v1d) -- [bend right=45, edge label'=$G_F$] (v2d),
  };
  \filldraw[black] (v1d) circle (2pt);
  \filldraw[black] (v2d) circle (2pt);
\end{feynman}
\node[font=\small] at (12.25, -12) {日落图（3 内线）};

% ===== Arrows =====
\draw[->, thick, green!70!black] (5.5, -4) -- (6.5, -4);
\draw[->, thick, green!70!black] (5.5, -7) -- (6.5, -7);
\draw[->, thick, green!70!black] (5.5, -10.5) -- (6.5, -10.5);
\node[font=\small\bfseries, green!70!black] at (6, -3.5) {Wick};

% ===== Legend =====
\node[font=\small, text width=6cm, align=left] at (13, -2) {%
  \textbf{读法}：\\
  $\bullet$ 黑点 = 相互作用顶点\\
  --- 实线 = 传播子 $G_F(x-y)$\\
  外线 = 未收缩场（外粒子）\\
  圈 = 圈积分 $\int\frac{d^4k}{(2\pi)^4}$
};

\end{tikzpicture}
\end{document}
